\section{Summary}

% --- SLIDE 1: SUMMARY ---
\begin{frame}[allowframebreaks]{Summary of the Lecture}

    \begin{itemize}
        \item \textbf{The ML code is the small box.} Most of the cost and most of the risk in a deployed
              model live in the infrastructure around it (Sculley et al., 2015).

        \item \textbf{The lifecycle is a loop}, not a line: experiment $\rightarrow$ validate $\rightarrow$
              package $\rightarrow$ deploy $\rightarrow$ monitor $\rightarrow$ retrain. MLOps exists to make
              that loop cheap enough to run often.

        \item \textbf{Reproducibility} comes from seeds, pinned environments, versioned data and declared
              configuration --- and it is what makes rollback possible.

        \item \textbf{Experiment tracking} turns a run into a first-class, append-only object with
              parameters, metrics and artifacts. MLflow: Tracking, Projects, Models, Model Registry.

        \item \textbf{Hyperparameter optimisation} beyond grid and random means \emph{using} past trials
              (TPE) and \emph{stopping} bad ones (successive halving / ASHA). Pruning is usually the bigger win.

        \item \textbf{Serving} is a choice between batch, online and streaming --- default to batch.
              Serialise the whole pipeline, and never load an untrusted pickle.

        \item \textbf{Monitoring} distinguishes upstream breakage, data drift and concept drift. Labels lag,
              so watch inputs and prediction distributions in the meantime.

        \item \textbf{Testing, staged rollout and lineage} are what let you ship a change without betting the
              product on it.
    \end{itemize}

\end{frame}

% --- SLIDE 2: THE MINIMUM VIABLE SETUP ---
\begin{frame}{If You Do Only Eight Things}

    \centering
    \textit{A weekend of work that removes most of the pain, for any project, at any scale.}
    \vspace{0.18em}

    \begin{columns}[T]
        \begin{column}{0.48\textwidth}
            \begin{tcolorbox}[colback=green!5, colframe=green!50!black, title=\scriptsize This week]
                \scriptsize
                \begin{enumerate}
                    \item Set every seed, and log it.
                    \item Commit a \textbf{lockfile}; log its hash with each run.
                    \item Write data to \textbf{immutable dated paths}; never overwrite.
                    \item Start a local \textbf{MLflow} server and log every run, including the failures.
                \end{enumerate}
            \end{tcolorbox}
        \end{column}
        \begin{column}{0.48\textwidth}
            \begin{tcolorbox}[colback=blue!5, colframe=blue!50, title=\scriptsize This month]
                \scriptsize
                \begin{enumerate}
                    \setcounter{enumi}{4}
                    \item Add a \textbf{schema check} at training \emph{and} serving input.
                    \item Serialise the \textbf{whole pipeline} as one artifact and register it by version.
                    \item \textbf{Log every prediction} with its model version.
                    \item Add \textbf{one} alert: prediction distribution versus a fixed reference window.
                \end{enumerate}
            \end{tcolorbox}
        \end{column}
    \end{columns}

    \vspace{0.24em}
    \begin{block}{The point}
        None of this requires a platform team, a cloud budget, or a new job title. It requires deciding that
        the run, the data and the model are \textbf{artifacts with identities} --- and everything else follows.
    \end{block}

\end{frame}

% --- SLIDE 3: PRACTICE / NEXT ---
\begin{frame}{Where to Practise This}

    \begin{columns}[T]
        \begin{column}{0.48\textwidth}
            \begin{tcolorbox}[colback=gray!8, colframe=gray!60, title=\small Exercise 1 --- Tracking and tuning]
                \small
                Wrap an existing training notebook in \textbf{MLflow} tracking; drive it with an
                \textbf{Optuna} study using a pruner; compare the runs in the UI; register the best model.
                \vspace{0.3em}
                \\ \textit{Entirely local --- no accounts, no cloud, no cost.}
            \end{tcolorbox}
        \end{column}
        \begin{column}{0.48\textwidth}
            \begin{tcolorbox}[colback=gray!8, colframe=gray!60, title=\small Exercise 2 --- Serving and drift]
                \small
                Export the trained pipeline, serve it with \textbf{FastAPI}, then simulate drift by shifting
                and corrupting the inputs; build a simple drift detector and alert on the shift in the
                prediction distribution.
            \end{tcolorbox}
        \end{column}
    \end{columns}

    \vspace{0.6em}
    \begin{itemize}
        \item \textbf{Then apply it to your own project.} Score yourself with the \textbf{ML Test Score}
              rubric and fix whichever of the four sections is lowest.
        \item \textbf{Where this continues:} governance, model cards and fairness auditing were covered
              earlier in this course. The training-loop reproducibility mechanics are in the
              computer-vision course; inference optimisation and distributed serving, for models too
              large for a single CPU, in the NLP course.
    \end{itemize}

\end{frame}
